\chapter{Introduction}\label{C:intro}
Recent development in the Internet of Things (IoT) have driven a
rapid increase the number of devices connected to the Internet, with
growth continues to accelerate in many areas: health care,
construction, building.
These devices typically communicate wirelessly, exchanging data over
radio frequency links instead of relying on wired connections, which
allows flexible deployment across wide range of physical environments.
% TODO: may be cite the growth of IoT
However, many such deploymens invole large number of constrained
devices, which are bounded by low computational power, memory and energy
(battery), to cover large geographic areas.
% TODO: what is LLNs
Because the
underlying network technologies are limited in range, multi-hop
routing is often required. Such networks are categorised as
Low-power and Lossy Networks (LLN) as links in LLN are
susceptible to loss and frequent packet drops in addtion to contrains
on individual devices.

These limitations led the Internet Engineering Task Force (IETF),
which is an international organisation of network designers and
researcher, to form Routing over Low-power and Lossy networks (ROLL)
working group to investigate the routing solution for LLNs.
The working group first published RFC 5548 \cite{rfc5548}, RFC 5673
\cite{rfc5673}, RFC 5826 \cite{rfc5826}, RFC 5867 \cite{rfc5867} to
define the routing requirements for smart urban, industrial area,
home automation and building automation applications respectively.
From there, the ROLL working group had investigated the existing routing
protocols Open Shortest Path First (OSPF), Intermediate System to
Intermediate System (IS-IS), Ad hoc On-Demand Distance Vector (AODV),
and Optimized Link State Routing Protocol (OLSR), to which they found out that
they do not satisfy all of
the routing requirements. % TODO: include the reason why these protocols faile
This led to ROLL introducing a new routing protocol
designed specifically for LLNs called IPv6 based routing protocol for
low power and lossy networks (RPL), standardised in RFC 6550
\cite{alexanderRPLIPv6Routing2012}.

At its core, RPL constructs a Destination-Oriented Directed Acyclic
Graph (DODAG) to route packets from nodes toward a root, using
objective function (OF) to calculate path cost and select parents
during DODAG construction. Because of this, OFs determine how the
traffic is distributed across DODAG
% TODO: and directly affect the link quality and power usage.
% TODO: talk a bit more about this

% TODO: make this better
The protocol fell short when it comes to performing under high traffic as
load balancing was not considered in the design. Furthermore, the OFs
that ROLL propose only consider one metric, which might not be enough
to capture the inherent conditions of the network
Moreover, most of
the OFs are reactive in nature, meaning that would react after a
route has degraded, leading to decrease in network quality and reliability.

% TODO: make this better
In this project, we propose a new machine learning based OF that take
load balancing into account and predict the link quality to choose the parents

\section{Contribution}
% TODO: Add one more contribution
\begin{itemize}
    \item Contribute a new objective function based on machine
        learning which focuses on improving delivery ratio
    \item Investigate the feature importance of different metrics in
        determine the delivery ratio quality
\end{itemize}

\section{Organisation}
The rest of this report is structured as follow
\begin{itemize}
    \item Chapter \ref{C:Background} provides background information
        in relation to RPL and objective functions and related work
        in improving the objective functions
    \item Chapter \ref{C:ProposedOf} provides the motivation and overview
        the machine learning techniques and models used for our
        proposed objective function
    \item Chapter \ref{C:Experiments} shows the experiment and result
        from applying ML to OF
\end{itemize}
